\documentclass[a4paper, 11pt]{article}
\usepackage[utf8]{inputenc}
\usepackage[T1]{fontenc}
\usepackage[margin=1.5cm]{geometry}

\title{Konspekt pracy inżynierskiej}
\author{Szymon Wasążnik\\325522}
\date{}

\begin{document}
\maketitle
\section{Cel}
Celem pracy jest opracowanie biblioteki w języku Python,
która umożliwi automatyczne optymalizowanie sieci neuronowej
pod względem architektury sieci na podstawie danych treningowych i walidacyjnych.
Optymalizowanie będzie odbywać się dzięki algorytmowi genetycznemu.

\section{Opis problematyki}
Dobór architektury sieci neuronowej i parametrów uczenia jest procesem,
wymagającym doświadczenia oraz licznych eksperymentów.
W praktyce trzeba testować wiele konfiguracji warstw,
liczby neuronów, funkcji aktywacji czy parametrów uczenia.
Do tego najlepiej testować daną konfigurację kilka razy,
ponieważ są losowane startowe wagi,
które mogą wpływać znacząco na końcowy wygląd sieci po uczeniu.

\noindent\\
Algorytmy genetyczne, pozwalają efektywnie przeszukiwać przestrzeń możliwych modeli
dzieki budowaniu ich na podstawie "genów".
Dzięki tej charakterystyce powinien dobrze sprawdzić się w dobieraniu odpowiedniej architektury i parametrów.

\noindent\\
Celem pracy jest zbudowanie biblioteki,
która automatyzuje proces dobierania przez wspomniany algorytm genetyczny.


\section{Zakres}
Optymalizacja odbywać się będzie względem podanych przez użytkownika danych testowych pod postacią tablic biblioteki \texttt{numpy}.


\section{Oczekiwane rezultaty}
Biblioteka w języku Python umożliwiająca wygenerowanie optymalnej architektury
sieci neuronowej.\\
Powinna ona dostarczać możliwie najoptymalniejsze rozwiązania bez potrzebnej ingerencji użytkownika.

\noindent\\
Dodatkowo 
co za tym idzie podczas projektu musi zostać znaleziona
\end{document}